% Notation macros for mini-unison.
%
% DELIBERATELY DISJOINT from macros.tex. arbor's \Store, \Names, \Cache, \nRef
% and friends denote arbor's objects, whose TYPES differ here: a Unison store is
% Hash -> Component, not Hash -> Node, and a Unison namespace is a relation, not
% a partial function. Overloading the notation would hide exactly the differences
% this document exists to state, so every object gets a \U-prefixed name and
% renders with a subscript u.

% --- the three stored structures ---
\newcommand{\UStore}{\Sigma_{\mathsf{u}}}      % codebase: Hash -> Component
\newcommand{\UNames}{N_{\mathsf{u}}}           % namespace: Name -> P(Referent)
\newcommand{\UCache}{W}                        % watch cache: Hash -> Hash
\newcommand{\UConfig}[3]{\langle #1,\, #2,\, #3\rangle}

% --- references and components ---
\newcommand{\Pos}{\mathsf{Pos}}
\newcommand{\Comp}{\mathsf{Comp}}
\newcommand{\Referent}{\mathsf{Referent}}
\newcommand{\Reference}{\mathsf{Reference}}
\newcommand{\UDer}[2]{\mathsf{Derived}(#1,#2)}
\newcommand{\UBuiltin}[1]{\mathsf{Builtin}(#1)}
\newcommand{\URefT}[1]{\mathsf{Ref}(#1)}
\newcommand{\UCon}[2]{\mathsf{Con}(#1,#2)}

% --- object language (named ABT, NOT de Bruijn) ---
\newcommand{\uvar}[1]{#1}
\newcommand{\ulam}[2]{\lambda #1.\,#2}
\newcommand{\uapp}[2]{#1\;#2}
\newcommand{\unat}[1]{\underline{#1}}
\newcommand{\utext}[1]{\text{``}#1\text{''}}
\newcommand{\uref}[1]{\mathsf{r}\,#1}
\newcommand{\uann}[2]{#1 : #2}

% --- types ---
\newcommand{\TNat}{\mathsf{Nat}}
\newcommand{\TText}{\mathsf{Text}}
\newcommand{\uarr}[2]{#1 \mathbin{\to} #2}
\newcommand{\utypes}[3]{#1 \vdash #2 : #3}       % Gamma |- t : T

% --- hashing ---
\newcommand{\uhash}[2]{\lceil #2 \rceil_{#1}}    % hash in environment
\newcommand{\uhashc}[1]{\lceil #1 \rceil}        % closed hash
\newcommand{\HEnv}{\rho}                          % the de-Bruijn-at-hash-time environment
\newcommand{\phash}[1]{\mathsf{ph}(#1)}           % permutation hash
\newcommand{\nhash}[1]{\mathsf{nh}(#1)}           % name hash (second sort key)
\newcommand{\comphash}{\mathsf{hashComp}}
\newcommand{\hashcomps}{\mathsf{hashComps}}
\newcommand{\amb}{\mathsf{Ambiguous}}
\newcommand{\bg}{\mathcal{B}}                     % a binding group

% --- naming ---
\newcommand{\unconf}[2]{\mathsf{unconflicted}(#1,#2)}
\newcommand{\uresolve}[3]{#1 \vdash #2 \rightsquigarrow #3}

% --- update ---
\newcommand{\deps}{\mathsf{deps}}
\newcommand{\dependents}{\mathsf{dependents}}
\newcommand{\tdeps}{\mathsf{tdeps}}
\newcommand{\ppe}{\mathsf{ppe}}
\newcommand{\render}{\mathsf{render}}
\newcommand{\uparse}{\mathsf{parse}}
\newcommand{\uupdate}{\mathsf{update}}
\newcommand{\lib}{\mathsf{lib}}
\newcommand{\scratch}{\mathsf{scratch}}
\newcommand{\ok}{\mathsf{typechecks}}
\newcommand{\usteps}{\longrightarrow_{\mathsf{u}}}
